\section{Appendix: Deployment Architecture Details}

This appendix is the wiring diagram stuff we pulled out of the main thread so reviewers can find RPC and ledger detail without mixing it into the RL theorem line.

\subsection{Layered Deployment Stack}

If you squint, the deployment stack is six slabs:
\begin{enumerate}
    \item \textbf{Network/RPC:} peer registration, transaction intake, and block propagation.
    \item \textbf{Blockchain core:} validation, mempool updates, and PoW checks.
    \item \textbf{PQC signatures:} post-quantum signing and verification interfaces.
    \item \textbf{Governance policy:} state encoding and action proposal.
    \item \textbf{Quantum uncertainty gate:} compact uncertainty estimate for fallback routing.
    \item \textbf{Shield + audit:} action validation and per-epoch commitment of audit records.
\end{enumerate}

\subsection{Why On-Chain Audit Commitments}

Participants do not trust each other's laptops. A local SQLite log is easy to ``lose.'' A hash on a shared ledger is a boring but useful witness when two orgs argue about who approved a knob turn.

\subsection{Boundary to Learning Claims}

Nothing in the appendix replaces the gated-and-shielded path evaluated in the tables. Chain commits are instrumentation for accountability; they do not move reward curves by themselves.
